\documentclass{article}



% Language setting

% Replace `english' with e.g. `spanish' to change the document language

\usepackage[english]{babel}



% Set page size and margins

% Replace `letterpaper' with `a4paper' for UK/EU standard size

\usepackage[letterpaper,top=2cm,bottom=2cm,left=3cm,right=3cm,marginparwidth=1.75cm]{geometry}



% Useful packages

\usepackage{amsmath}

\usepackage{amssymb}

\usepackage{graphicx}

\usepackage[colorlinks=true, allcolors=black]{hyperref}

\usepackage{titlesec}
\usepackage{xparse}



\NewDocumentCommand{\qcq}{m o o m}{%
	\ensuremath{\forall #1 \IfValueT{#2}{, #2} \IfValueT{#3}{, #3}  \in  \mathbb{#4}}%
}
\NewDocumentCommand{\ext}{m o o m}{%
	\ensuremath{\exists #1 \IfValueT{#2}{, #2} \IfValueT{#3}{, #3}  \in  \mathbb{#4}}%
}



\title{\textbf{Limites}}

\author{\textbf{Yassin Akkab}}



\newcommand{\R}{\mathbb{R}}
\newcommand{\N}{\mathbb{N}}
\newcommand{\Z}{\mathbb{Z}}
\newcommand{\C}{\mathbb{C}}
\newcommand{\Q}{\mathbb{Q}}
\newcommand{\D}{\mathbb{D}}





\graphicspath{ {./Screenshots/} }




\begin{document}
	
	\maketitle
	
	
	
	\tableofcontents
	
	\newpage

	
	\section{Limite d’une fonction}
	\subsection{Définition}
		Soient $f : I \to \mathbb{R}$, $a \in \mathbb{R}$ adhérent à $I$ et $l \in \mathbb{R}$. Alors on définit le fait que $f$ tende vers $l$ en $a$ de l’une des 9 manières suivantes :
		\begin{enumerate}
			\item Si $a \in \mathbb{R}$ :
			\begin{enumerate}
				\item On dit que $f$ tend vers $l \in \mathbb{R}$ lorsque $x$ tend vers $a$, ce que l’on note $\lim\limits_{x \to a} f(x) = l$, si :
				\[
				\forall \epsilon > 0, \exists \eta > 0, \forall x \in I, |x - a| \leq \eta \Rightarrow |f(x) - l| \leq \epsilon;
				\]
				\item On dit que $f$ tend vers $+\infty$ en $a$, ce que l’on note $\lim\limits_{x \to a} f(x) = +\infty$, si :
				\[
				\forall A \in \mathbb{R}, \exists \eta > 0, \forall x \in I, |x - a| \leq \eta \Rightarrow f(x) \geq A;
				\]
				\item On dit que $f$ tend vers $-\infty$ en $a$, ce que l’on note $\lim\limits_{x \to a} f(x) = -\infty$, si :
				\[
				\forall A \in \mathbb{R}, \exists \eta > 0, \forall x \in I, |x - a| \leq \eta \Rightarrow f(x) \leq A;
				\]
			\end{enumerate}
			
			\item Si $a = +\infty$ :
			\begin{enumerate}
				\item On dit que $f$ tend vers $l \in \mathbb{R}$ lorsque $x$ tend vers $+\infty$, ce que l’on note $\lim\limits_{x \to +\infty} f(x) = l$, si :
				\[
				\forall \epsilon > 0, \exists B \in \mathbb{R}, \forall x \in I, x \geq B \Rightarrow |f(x) - l| \leq \epsilon;
				\]
				\item On dit que $f$ tend vers $+\infty$ en $+\infty$, ce que l’on note $\lim\limits_{x \to +\infty} f(x) = +\infty$, si :
				\[
				\forall A \in \mathbb{R}, \exists B \in \mathbb{R}, \forall x \in I, x \geq B \Rightarrow f(x) \geq A;
				\]
				\item On dit que $f$ tend vers $-\infty$ en $+\infty$, ce que l’on note $\lim\limits_{x \to +\infty} f(x) = -\infty$, si :
				\[
				\forall A \in \mathbb{R}, \exists B \in \mathbb{R}, \forall x \in I, x \geq B \Rightarrow f(x) \leq A;
				\]
			\end{enumerate}
			
			\item Si $a = -\infty$ :
			\begin{enumerate}
				\item On dit que $f$ tend vers $l \in \mathbb{R}$ lorsque $x$ tend vers $-\infty$, ce que l’on note $\lim\limits_{x \to -\infty} f(x) = l$, si :
				\[
				\forall \epsilon > 0, \exists B \in \mathbb{R}, \forall x \in I, x \leq B \Rightarrow |f(x) - l| \leq \epsilon;
				\]
				\item On dit que $f$ tend vers $+\infty$ en $-\infty$, ce que l’on note $\lim\limits_{x \to -\infty} f(x) = +\infty$, si :
				\[
				\forall A \in \mathbb{R}, \exists B \in \mathbb{R}, \forall x \in I, x \leq B \Rightarrow f(x) \geq A;
				\]
				\item On dit que $f$ tend vers $-\infty$ en $-\infty$, ce que l’on note $\lim\limits_{x \to -\infty} f(x) = -\infty$, si :
				\[
				\forall A \in \mathbb{R}, \exists B \in \mathbb{R}, \forall x \in I, x \leq B \Rightarrow f(x) \leq A;
				\]
			\end{enumerate}
		\end{enumerate}
	\subsection{Proposition (Unicité de la limite)}
	
	Si une fonction tend vers une limite en $a$, alors cette limite est unique.
	
	\subsubsection*{Démonstration}
	
	Montrons le par exemple dans le cas où $a \in \mathbb{R}$ avec des limites finies.
	
	Par l'absurde, supposons que la fonction $f$ définie sur $I$ tende simultanément en $a$ vers $l_1, l_2 \in \mathbb{R}$, avec $l_1 \neq l_2$. Posons $\varepsilon = \frac{|l_1 - l_2|}{3}$. Par définition de la limite de $f$ en $a$ :
	
	\begin{itemize}
		\item Il existe $\eta_1 > 0$ tel que, pour tout $x \in I$, si $|x - a| \leq \eta_1$, alors $|f(x) - l_1| \leq \varepsilon$.
		\item Il existe $\eta_2 > 0$ tel que, pour tout $x \in I$, si $|x - a| \leq \eta_2$, alors $|f(x) - l_2| \leq \varepsilon$.
	\end{itemize}
	
	En posant $\eta = \min(\eta_1, \eta_2)$, considérons un $x \in I$ tel que $|x - a| \leq \eta$ (ce qui est possible, puisque $a$ est adhérent à $I$). On a alors, pour un tel $x$ :
	
	\[
	3\varepsilon = |l_1 - l_2| = |f(x) - l_1 + l_2 - f(x)| \leq |f(x) - l_1| + |f(x) - l_2| \leq 2\varepsilon,
	\]
	
	d'où la contradiction avec le fait que $\varepsilon > 0$.
\section{Limites latérales}

Les limites latérales permettent d'étudier le comportement d'une fonction lorsque la variable approche une certaine valeur \( a \), soit par la gauche, soit par la droite. Elles sont particulièrement utiles pour étudier les discontinuités d'une fonction.

\subsection{Définitions}

\begin{itemize}
	\item La \textbf{limite à gauche} de \( f(x) \) lorsque \( x \) tend vers \( a \) est notée :
	\[
	\lim_{x \to a^-} f(x)
	\]
	Cela signifie que \( x \) approche \( a \) par des valeurs inférieures à \( a \) (par la gauche).
	
	\item La \textbf{limite à droite} de \( f(x) \) lorsque \( x \) tend vers \( a \) est notée :
	\[
	\lim_{x \to a^+} f(x)
	\]
	Cela signifie que \( x \) approche \( a \) par des valeurs supérieures à \( a \) (par la droite).
	
	\item Si les deux limites latérales sont égales, alors la limite en \( a \) existe et est égale à cette valeur :
	\[
	\lim_{x \to a} f(x) = L \quad \text{si et seulement si} \quad \lim_{x \to a^-} f(x) = \lim_{x \to a^+} f(x) = L
	\]
	Sinon, il y a une discontinuité en \( a \).
\end{itemize}

\subsection{Exemple de limites latérales}

Considérons la fonction définie par morceaux suivante :
\[
f(x) =
\begin{cases}
	x + 2 & \text{si } x < 1 \\
	3x - 1 & \text{si } x \geq 1
\end{cases}
\]

Calculons les limites latérales en \( x = 1 \) :
\begin{itemize}
	\item Limite à gauche :
	\[
	\lim_{x \to 1^-} f(x) = \lim_{x \to 1^-} (x + 2) = 1 + 2 = 3
	\]
	\item Limite à droite :
	\[
	\lim_{x \to 1^+} f(x) = \lim_{x \to 1^+} (3x - 1) = 3(1) - 1 = 2
	\]
\end{itemize}

Dans cet exemple, comme \( \lim_{x \to 1^-} f(x) \neq \lim_{x \to 1^+} f(x) \), la limite en \( x = 1 \) n'existe pas et il y a une discontinuité en \( x = 1 \).

\subsection{Utilisation des limites latérales}

Les limites latérales sont souvent utilisées pour :
\begin{itemize}
	\item Étudier le comportement des fonctions définies par morceaux.
	\item Déterminer la continuité d'une fonction en un point donné.
	\item Identifier les discontinuités de saut ou les discontinuités infinies.
\end{itemize}


	

	
	\subsubsection{Remarque}
	
	On peut aussi définir les limites à gauche et à droite avec des $\varepsilon$. Par exemple, dire que $f$ a pour limite à gauche $l$ en $a$ se traduit par :
	
	\begin{itemize}
		\item Si $l \in \mathbb{R}$ : 
		\[
		\forall \varepsilon > 0, \exists \eta > 0, \forall x \in I, a - \eta \leq x < a \Rightarrow |f(x) - l| \leq \varepsilon;
		\]
		\item Si $l = +\infty$ :
		\[
		\forall A \in \mathbb{R}, \exists \eta > 0, \forall x \in I, a - \eta \leq x < a \Rightarrow f(x) \geq A;
		\]
		\item Si $l = -\infty$ :
		\[
		\forall A \in \mathbb{R}, \exists \eta > 0, \forall x \in I, a - \eta \leq x < a \Rightarrow f(x) \leq A.
		\]
	\end{itemize}
	\subsubsection{Proposition}
	
	Si $a \in \mathbb{R}$, $l \in \mathbb{R}$, et $f$ est définie sur $I$, à gauche et à droite de $a$. Alors :
	
	\begin{enumerate}
		\item Si $a \in I$ : 
		\[
		\lim_{x \to a} f(x) = l \quad \text{si, et seulement si,} \quad \lim_{x \to a^+} f(x) = \lim_{x \to a^-} f(x) = f(a) = l.
		\]
		
		\item Si $a \notin I$ : 
		\[
		\lim_{x \to a} f(x) = l \quad \text{si, et seulement si,} \quad \lim_{x \to a^+} f(x) = \lim_{x \to a^-} f(x) = l.
		\]
	\end{enumerate}
	

	
	\subsubsection{Exemples}
	
	\begin{enumerate}
		\item On peut voir ainsi que la fonction inverse n’a pas de limite en 0, puisque :
		\[
		\lim_{x \to 0^+} \frac{1}{x} = +\infty \quad \text{et} \quad \lim_{x \to 0^-} \frac{1}{x} = -\infty.
		\]
		
		\item En revanche, la fonction $x \mapsto \frac{1}{x^2}$ a bien une limite en 0 puisque :
		\[
		\lim_{x \to 0^+} \frac{1}{x^2} = +\infty = \lim_{x \to 0^-} \frac{1}{x^2}.
		\]
	\end{enumerate}
	
	\section{Somme, produit et rapport avec des constantes et des infinis}
	
	Dans l'étude des limites, il est important de connaître les règles qui s'appliquent lorsqu'on combine des constantes, des infinis, et des fonctions dans des opérations telles que la somme, le produit ou le rapport. Voici quelques résultats clés :
	
	\subsection{Somme avec des infinis et des constantes}
	\begin{itemize}
		\item $\infty + \infty = \infty$
		\item $-\infty + (-\infty) = -\infty$
		\item $\infty + c = \infty$ \quad pour toute constante réelle \( c \)
		\item $-\infty + c = -\infty$ \quad pour toute constante réelle \( c \)
		\item $\infty + (-\infty)$ est une \textit{forme indéterminée}.
	\end{itemize}
	
	\subsection{Produit avec des infinis et des constantes}
	\begin{itemize}
		\item $c \cdot \infty = \infty$ \quad pour \( c > 0 \)
		\item $c \cdot \infty = -\infty$ \quad pour \( c < 0 \)
		\item $c \cdot 0 = 0$ \quad pour toute constante \( c \)
		\item $\infty \cdot \infty = \infty$
		\item $\infty \cdot 0$ est une \textit{forme indéterminée}.
	\end{itemize}
	
	\subsection{Rapport avec des infinis et des constantes}
	\begin{itemize}
		\item $\frac{\infty}{c} = \infty$ \quad pour \( c > 0 \)
		\item $\frac{-\infty}{c} = -\infty$ \quad pour \( c > 0 \)
		\item $\frac{\infty}{\infty}$ est une \textit{forme indéterminée}.
		\item $\frac{c}{\infty} = 0$ \quad pour toute constante réelle \( c \)
		\item $\frac{c}{0}$ pour \( c \neq 0 \) tend vers \( \pm \infty \) selon le signe de \( c \).
	\end{itemize}
	
	Ces règles sont essentielles pour comprendre les comportements asymptotiques des fonctions, notamment dans les calculs de limites et dans l'étude des séries et des intégrales impropres.
	
	\section{Techniques de calcul des limites indéterminées}
	
	\subsection{Les quatre formes indéterminées des limites}
	
	Lors de l'étude des limites, il existe quatre formes indéterminées classiques :
	
	\begin{itemize}
		\item $\frac{0}{0}$ : Cette forme apparaît souvent dans les limites de fractions lorsque le numérateur et le dénominateur tendent simultanément vers zéro.
		\item $\frac{\infty}{\infty}$ : Cette forme survient lorsque le numérateur et le dénominateur d'une fraction tendent tous deux vers l'infini.
		\item $0 \times \infty$ : Ce produit indéterminé peut se rencontrer dans des expressions où un terme tend vers zéro tandis qu'un autre tend vers l'infini.
		\item $\infty - \infty$ : Cette différence indéterminée apparaît lorsque deux termes tendent vers l'infini, mais leur différence n'est pas immédiatement claire.
	\end{itemize}
	
	Ces formes nécessitent des techniques spéciales pour être résolues, telles que la factorisation, la rationalisation, l'encadrement, ou l'utilisation des théorèmes comme celui de l'Hôpital.
	
	\subsection{Proposition }
	Si \( f \), \( g \) sont deux fonctions définies sur \( I \). Si \( f(x) \leq g(x) \) alors :
	

	\begin{enumerate}
		\item  	Si \( f \) et \( g \) ont des limites en \( a \)	\(\lim_{x \to a} f(x) \leq \lim_{x \to a} g(x) \);
		\item si \( \lim_{x \to a} f(x) = +\infty \), alors \( \lim_{x \to a} g(x) = +\infty \);
		\item si \( \lim_{x \to a} g(x) = -\infty \), alors \( \lim_{x \to a} f(x) = -\infty \).
	\end{enumerate}
	\subsection{Théorème  (Théorème d'encadrement, ou des gendarmes)}
	Si \( f \), \( g \), \( h \) sont trois fonctions définies sur \( I \), \( a \) adhérent à \( I \), tels qu'au voisinage de \( a \) : \( f(x) \leq g(x) \leq h(x) \). Si
	\[
	\lim_{x \to a} f(x) = \lim_{x \to a} h(x) = l \in \mathbb{R},
	\]
	alors
	\[
	\lim_{x \to a} g(x) = l.
	\]
	\subsection{Le théorème de l'Hôpital}
	
	Le théorème de l'Hôpital est un outil puissant pour résoudre les limites présentant des formes indéterminées, telles que $\frac{0}{0}$ et $\frac{\infty}{\infty}$. Il permet de simplifier ces formes en dérivant les numérateurs et dénominateurs sous certaines conditions.
	\\
	
	\textbf{Théorème de l'Hôpital} : \textit{Soient $f$ et $g$ deux fonctions dérivables sur un intervalle ouvert $I$ contenant un point $a$ (à l'exception peut-être de $a$ lui-même). Si :
		\begin{itemize}
			\item $\lim_{x \to a} f(x) = \lim_{x \to a} g(x) = 0$ ou $\lim_{x \to a} f(x) = \lim_{x \to a} g(x) = \pm \infty$,
			\item $g'(x) \neq 0$ pour tout $x$ proche de $a$,
		\end{itemize}
		alors :
		\[
		\lim_{x \to a} \frac{f(x)}{g(x)} = \lim_{x \to a} \frac{f'(x)}{g'(x)}
		\]
		si la limite à droite existe ou tend vers $\pm \infty$.}
	
	\textbf{Remarques} :
	\begin{itemize}
		\item Le théorème de l'Hôpital peut être appliqué plusieurs fois si nécessaire, notamment lorsque la forme indéterminée persiste après une première application.
		\item Il ne peut être appliqué qu'en présence de formes indéterminées spécifiques comme $\frac{0}{0}$ ou $\frac{\infty}{\infty}$.
	\end{itemize}
	
	Le théorème est souvent utilisé en combinaison avec d'autres techniques de calcul de limites pour obtenir un résultat plus simple et gérable.
	
\section{Limites usuelles}

\subsection{Limites des fonctions usuelles dans un point}
Soient \( f \) et \( g \) deux fonctions définies sur un intervalle \( I \subset \mathbb{R} \), et \( a \in I \). On a les résultats suivants pour les limites usuelles :

\begin{itemize}
	\item \textbf{Limite d'une constante :} Si \( f(x) = c \) où \( c \in \mathbb{R} \), alors :
	\[
	\lim_{x \to a} f(x) = c
	\]
	
	\item \textbf{Limite de l'identité :} Si \( f(x) = x \), alors :
	\[
	\lim_{x \to a} f(x) = a
	\]
	
	\item \textbf{Limite de l'inverse :} Si \( f(x) = \frac{1}{x} \) et \( a \neq 0 \), alors :
	\[
	\lim_{x \to a} \frac{1}{x} = \frac{1}{a}
	\]
	
	\item \textbf{Limite de la fonction carrée :} Si \( f(x) = x^2 \), alors :
	\[
	\lim_{x \to a} x^2 = a^2
	\]
	
	\item \textbf{Limite de la fonction racine carrée :} Si \( f(x) = \sqrt{x} \) et \( a \geq 0 \), alors :
	\[
	\lim_{x \to a} \sqrt{x} = \sqrt{a}
	\]
	
	

\end{itemize}

Ces résultats sont essentiels pour le calcul des limites en analyse et seront utilisés dans de nombreux cas.


\subsection{Proposition. Limites usuelles des fonctions en \( +\infty \) et \( -\infty \)}

\subsubsection{Limites des fonctions polynômes}
Soit \( f(x) = a_n x^n + a_{n-1} x^{n-1} + \cdots + a_1 x + a_0 \), un polynôme de degré \( n \). Alors :
\[
\lim_{x \to +\infty} f(x) =
\begin{cases} 
	+ \infty & \text{si } a_n > 0, \\
	- \infty & \text{si } a_n < 0, 
\end{cases}
\quad \text{et} \quad
\lim_{x \to -\infty} f(x) =
\begin{cases} 
	+ \infty & \text{si } n \text{ est pair et } a_n > 0, \\
	- \infty & \text{si } n \text{ est pair et } a_n < 0, \\
	+ \infty & \text{si } n \text{ est impair et } a_n < 0, \\
	- \infty & \text{si } n \text{ est impair et } a_n > 0.
\end{cases}
\]

\subsubsection{Limites des fonctions rationnelles}
Soient \( f(x) = \frac{p(x)}{q(x)} \), où \( p(x) \) et \( q(x) \) sont des polynômes tels que \( \deg(p) = m \) et \( \deg(q) = n \). Alors :
\[
\lim_{x \to +\infty} \frac{f(x)}{g(x)} =
\begin{cases} 
	0 & \text{si } m < n, \\
	\frac{a_m}{b_n} & \text{si } m = n, \\
	+ \infty & \text{si } m > n \text{ et } \frac{a_m}{b_n} > 0, \\
	- \infty & \text{si } m > n \text{ et } \frac{a_m}{b_n} < 0,
\end{cases}
\quad \text{et} \quad
\lim_{x \to -\infty} \frac{f(x)}{g(x)} = \text{à calculer en fonction des signes de } a_m \text{ et } b_n.
\]

\subsubsection{Limites des fonctions racines carrées}
Soit \( f(x) = \sqrt{x} \). Alors :
\[
\lim_{x \to +\infty} \sqrt{x} = +\infty,
\quad \lim_{x \to 0^+} \sqrt{x} = 0,
\quad \text{et} \quad \lim_{x \to 0^-} \sqrt{x} \text{ n'existe pas}.
\]



\subsubsection{Les limites des formes trigonométriques}

Les limites faisant intervenir des fonctions trigonométriques sont fréquemment rencontrées dans l'étude des fonctions. Voici quelques résultats importants à connaître :

\begin{itemize}
	\item $\lim_{x \to 0} \frac{\sin(x)}{x} = 1$
	\item $\lim_{x \to 0} \frac{1 - \cos(x)}{x^2} = \frac{1}{2}$
	\item $\lim_{x \to \infty} \sin(x)$ et $\lim_{x \to \infty} \cos(x)$ n'existent pas car ces fonctions oscillent entre -1 et 1.
	\item $\lim_{x \to 0} \tan(x) = 0$
	\item $\lim_{x \to 0} \frac{\tan(x)}{x} = 1$
	\item $\lim_{x \to \frac{\pi}{2}^-} \tan(x) = +\infty$
	\item $\lim_{x \to \frac{\pi}{2}^+} \tan(x) = -\infty$
	\item $\lim_{x \to -\frac{\pi}{2}^-} \tan(x) = -\infty$
	\item $\lim_{x \to -\frac{\pi}{2}^+} \tan(x) = +\infty$
\end{itemize}

Ces limites sont essentielles pour résoudre des problèmes complexes en analyse, notamment dans les développements en série et l'application du\textbf{ théorème de l'Hôpital}.


	

	

	
\end{document}